%% ============================================================================
%%  Lab Report – Agentic AI: Travel Planning with LangChain & MCP
%%  Compile:  pdflatex report.tex  (×2 for cross-refs)
%% ============================================================================
\documentclass[12pt,a4paper]{article}

% ── Packages ──────────────────────────────────────────────────────────────── %
\usepackage[utf8]{inputenc}
\usepackage[T1]{fontenc}
\usepackage[english]{babel}
\usepackage{geometry}
\geometry{margin=2.5cm}
\usepackage{graphicx}
\usepackage{xcolor}
\usepackage{hyperref}
\usepackage{listings}
\usepackage{booktabs}
\usepackage{amsmath}
\usepackage{caption}
\usepackage{float}
\usepackage{enumitem}
\usepackage{fancyhdr}
\usepackage{titlesec}
\usepackage{tcolorbox}
\tcbuselibrary{listingsutf8,skins}

% ── Colors ────────────────────────────────────────────────────────────────── %
\definecolor{codeblue}{HTML}{667eea}
\definecolor{codegray}{HTML}{6c7086}
\definecolor{codebg}{HTML}{f5f5f5}
\definecolor{linkcolor}{HTML}{764ba2}

\hypersetup{
  colorlinks=true,
  linkcolor=linkcolor,
  urlcolor=codeblue,
  citecolor=codegray,
}

% ── Code listing style ────────────────────────────────────────────────────── %
\lstdefinestyle{python}{
  language=Python,
  basicstyle=\ttfamily\small,
  backgroundcolor=\color{codebg},
  keywordstyle=\color{codeblue}\bfseries,
  commentstyle=\color{codegray}\itshape,
  stringstyle=\color{linkcolor},
  breaklines=true,
  frame=single,
  rulecolor=\color{codegray!30},
  showstringspaces=false,
  tabsize=4,
  xleftmargin=1em,
  framexleftmargin=0.5em,
}
\lstset{style=python}

% ── Header / Footer ──────────────────────────────────────────────────────── %
\pagestyle{fancy}
\fancyhf{}
\fancyhead[L]{\small\textit{Lab – Agentic AI}}
\fancyhead[R]{\small\textit{LangChain \& MCP}}
\fancyfoot[C]{\thepage}
\renewcommand{\headrulewidth}{0.4pt}

% ── Tip box ───────────────────────────────────────────────────────────────── %
\newtcolorbox{infobox}[1][]{
  colback=codeblue!5, colframe=codeblue!60,
  fonttitle=\bfseries, title={#1}, arc=3pt,
}

% ══════════════════════════════════════════════════════════════════════════════
\begin{document}

% ── Title page ────────────────────────────────────────────────────────────── %
\begin{titlepage}
\centering
\vspace*{3cm}
{\Huge\bfseries Lab: Agentic AI\\[0.4em]}
{\LARGE Travel Planning Assistant\\with LangChain, LangGraph \& MCP\par}
\vspace{1.5cm}
{\Large\itshape Yassine\par}
\vspace{1cm}
{\large University / Institution Name\par}
\vspace{0.5cm}
{\large Academic Year 2025--2026\par}
\vfill
{\large \today\par}
\end{titlepage}

\begin{center}
\fcolorbox{codeblue}{codeblue!5}{%
  \parbox{0.85\textwidth}{\centering\vspace{0.4em}
    \textbf{GitHub Repository:}\\[0.3em]
    \large\url{https://github.com/Moad26/lab-agentic}
  \vspace{0.4em}}%
}
\end{center}
\vspace{1em}

\tableofcontents
\newpage

% ══════════════════════════════════════════════════════════════════════════════
\section{Introduction}
% ══════════════════════════════════════════════════════════════════════════════

Large Language Models (LLMs) excel at generating fluent text, but they cannot
natively access external data or perform precise computations.
\textbf{Agentic AI} bridges this gap by embedding an LLM inside a
\emph{reason-then-act} loop: the model observes a user query, decides which
external \textbf{tool} to invoke, interprets the result, and repeats until it
can deliver a final answer.

In this lab we build an \textbf{Agentic Travel Planning Assistant} that
combines:
\begin{itemize}[nosep]
  \item \textbf{Ollama} – local inference of open-weight models (LLaMA~3.1);
  \item \textbf{LangChain / LangGraph} – the orchestration framework that
        implements a \emph{ReAct} agent graph;
  \item \textbf{Model Context Protocol (MCP)} – a standardised protocol for
        exposing tools as lightweight HTTP servers;
  \item \textbf{Streamlit} – a web UI that surfaces tool calls, critic
        reviews, and budget constraints.
\end{itemize}

Three exercises progressively enhance the baseline agent:
\begin{enumerate}[nosep]
  \item \textbf{Tool-call transparency} – every MCP call is logged and
        rendered in the GUI.
  \item \textbf{Critic agent} – a second LLM pass reviews the itinerary for
        feasibility, budget, and balance.
  \item \textbf{Budget constraint} – the user sets a maximum budget; the agent
        must respect it, with one-click re-planning.
\end{enumerate}

% ══════════════════════════════════════════════════════════════════════════════
\section{Background}
% ══════════════════════════════════════════════════════════════════════════════

\subsection{ReAct Paradigm}

The \emph{Reasoning + Acting} (ReAct) framework~\cite{yao2023react} interleaves
chain-of-thought reasoning with tool actions.  At each step the agent produces:
\begin{enumerate}[nosep]
  \item \textbf{Thought} – internal reasoning about the current state;
  \item \textbf{Action} – a tool call with structured arguments;
  \item \textbf{Observation} – the tool's response, fed back as context.
\end{enumerate}
This loop continues until the agent decides to emit a \texttt{Final Answer}.

\subsection{Model Context Protocol (MCP)}

MCP is an open protocol, introduced by Anthropic in late 2024, that
standardises how LLM applications discover and invoke tools.  Each
\textbf{MCP server} exposes one or more tools via Server-Sent Events (SSE)
or \texttt{stdio} transport.  The client (here, \texttt{langchain-mcp-adapters})
connects to the servers, retrieves tool schemas, and converts them into
LangChain \texttt{Tool} objects automatically.

\subsection{LangGraph}

LangGraph models agent execution as a directed graph of \emph{nodes}
(LLM calls, tool dispatchers) and \emph{edges} (conditional routing).
The built-in \texttt{create\_react\_agent()} helper constructs a complete
ReAct graph with a tool node that dispatches calls to the appropriate
LangChain tool.

% ══════════════════════════════════════════════════════════════════════════════
\section{Project Architecture}
% ══════════════════════════════════════════════════════════════════════════════

\subsection{Repository Layout}

After restructuring, the project follows a clean \texttt{src/} layout:

\begin{verbatim}
lab-agentic/
├── src/
│   ├── agent/
│   │   ├── __init__.py
│   │   ├── agent.py          # ReAct agent + critic
│   │   └── app.py            # Streamlit web UI
│   └── mcp_servers/
│       ├── travel_search_server.py   # port 3001
│       ├── finance_server.py         # port 3002
│       ├── weather_server.py         # port 3003
│       ├── currency_server.py        # port 3004
│       └── calculator_server.py      # port 3005
├── scripts/
│   └── run_servers.py
├── tests/
├── docs/
│   └── report.tex
├── pyproject.toml
└── README.md
\end{verbatim}

\subsection{Component Diagram}

Figure~\ref{fig:arch} illustrates the runtime data flow.

\begin{figure}[H]
\centering
\begin{tcolorbox}[colback=white,colframe=codegray!50,width=0.95\textwidth,arc=4pt]
\begin{verbatim}
  User Query  ──►  Streamlit UI (app.py)
                        │
                        ▼
               LangGraph ReAct Agent (agent.py)
                  LLM: Ollama / llama3.1
                        │
            ┌───────────┼───────────┐
            ▼           ▼           ▼
      MCP Server 1  MCP Server 2  ...  MCP Server 5
      (Travel)      (Finance)          (Calculator)
      :3001         :3002              :3005
            │           │           │
            └───────────┼───────────┘
                        ▼
               Agent Final Answer
                        │
                        ▼  (optional)
               Critic Agent Review
\end{verbatim}
\end{tcolorbox}
\caption{Runtime architecture of the Agentic Travel Planner.}
\label{fig:arch}
\end{figure}

\subsection{Technology Stack}

\begin{table}[H]
\centering
\caption{Key dependencies and their roles.}
\label{tab:deps}
\begin{tabular}{@{}llp{7cm}@{}}
\toprule
\textbf{Library} & \textbf{Version} & \textbf{Role} \\
\midrule
\texttt{langchain}             & $\geq$ 0.3  & LLM orchestration, prompt templates \\
\texttt{langgraph}             & $\geq$ 0.4  & ReAct agent graph construction \\
\texttt{langchain-ollama}      & $\geq$ 0.3  & \texttt{ChatOllama} LLM wrapper \\
\texttt{langchain-mcp-adapters}& $\geq$ 0.1  & MCP $\leftrightarrow$ LangChain tool bridge \\
\texttt{mcp[cli]}              & $\geq$ 1.0  & \texttt{FastMCP} server framework \\
\texttt{streamlit}             & $\geq$ 1.40 & Web UI \\
\texttt{ollama}                & --           & Local LLM inference (llama3.1) \\
\bottomrule
\end{tabular}
\end{table}

% ══════════════════════════════════════════════════════════════════════════════
\section{MCP Tool Servers}
% ══════════════════════════════════════════════════════════════════════════════

Each MCP server is a standalone Python script built with \texttt{FastMCP}.
They expose tools via SSE on \texttt{localhost}.

\subsection{Travel Search Server (port 3001)}

Provides a simulated database of five destinations (Barcelona, Paris, Tokyo,
New York, Marrakech) with attractions and activities.

\begin{lstlisting}
@mcp.tool()
def search_destination(destination: str) -> str:
    """Search for attractions and activities
    for a given destination city."""
    key = destination.strip().lower()
    # partial match lookup ...
    return formatted_result

@mcp.tool()
def list_available_destinations() -> str:
    """List all destinations in the database."""
\end{lstlisting}

\subsection{Finance Server (port 3002)}

Estimates travel budgets using per-destination daily cost profiles
(accommodation, food, transport, activities, miscellaneous) and
round-trip flight estimates, all in USD.

\begin{lstlisting}
@mcp.tool()
def estimate_budget(destination: str, days: int) -> str:
    """Estimate total travel budget in USD."""

@mcp.tool()
def get_daily_cost(destination: str) -> str:
    """Get average daily cost breakdown."""
\end{lstlisting}

\subsection{Weather Server (port 3003)}

Returns monthly temperature ranges, conditions, rainy-day counts, and
indoor/outdoor recommendations.  A nearest-month fallback handles months
missing from the database.

\subsection{Currency Server (port 3004)}

Converts amounts between 15 currencies using static exchange rates against
USD.

\subsection{Calculator Server (port 3005)}

Evaluates safe arithmetic expressions, computes percentages, and splits
costs among travellers.

\begin{table}[H]
\centering
\caption{Summary of all MCP tools.}
\label{tab:tools}
\begin{tabular}{@{}lll@{}}
\toprule
\textbf{Server} & \textbf{Port} & \textbf{Tools} \\
\midrule
Travel Search & 3001 & \texttt{search\_destination}, \texttt{list\_available\_destinations} \\
Finance       & 3002 & \texttt{estimate\_budget}, \texttt{get\_daily\_cost} \\
Weather       & 3003 & \texttt{get\_weather}, \texttt{get\_best\_months} \\
Currency      & 3004 & \texttt{convert\_currency}, \texttt{list\_currencies} \\
Calculator    & 3005 & \texttt{calculate}, \texttt{percentage}, \texttt{split\_cost} \\
\bottomrule
\end{tabular}
\end{table}

% ══════════════════════════════════════════════════════════════════════════════
\section{Agent Implementation}
% ══════════════════════════════════════════════════════════════════════════════

\subsection{ReAct Agent (\texttt{agent.py})}

The main agent is constructed with a single function call:

\begin{lstlisting}
async def run_travel_agent(user_request, model="llama3.1",
                           budget_limit=None):
    llm = ChatOllama(model=model, temperature=0.3)

    async with MultiServerMCPClient(MCP_SERVERS) as client:
        tools = client.get_tools()
        agent = create_react_agent(llm, tools)
        result = await agent.ainvoke({
            "messages": [
                ("system", SYSTEM_PROMPT),
                ("user", prompt),
            ]
        })
    # parse messages -> extract tool_calls, output
    return {"output": ..., "tool_calls": [...], "messages": [...]}
\end{lstlisting}

\textbf{Key design choices:}
\begin{itemize}[nosep]
  \item \texttt{MultiServerMCPClient} connects to all five servers
        simultaneously and merges their tool schemas.
  \item The system prompt instructs the agent on the available tools and
        a six-step reasoning process.
  \item \texttt{temperature=0.3} balances creativity with consistency.
  \item Post-processing parses the message list to extract a structured
        log of tool calls (name, input, output).
\end{itemize}

\subsection{System Prompt Engineering}

The system prompt is critical to agent behavior.  It specifies:

\begin{enumerate}[nosep]
  \item The available tools grouped by category;
  \item A step-by-step workflow (understand $\to$ search $\to$ budget $\to$
        weather $\to$ currency $\to$ calculate $\to$ synthesise);
  \item The output format: a structured, day-by-day travel itinerary.
\end{enumerate}

% ══════════════════════════════════════════════════════════════════════════════
\section{Exercises}
% ══════════════════════════════════════════════════════════════════════════════

\subsection{Exercise 1 – Tool-Call Transparency}

\textbf{Goal:} Display every MCP tool call in the Streamlit UI so the user
can inspect the agent's reasoning process.

\textbf{Implementation:} After the agent completes, the message history is
iterated.  Each \texttt{AIMessage} with a \texttt{tool\_calls} attribute is
logged, and the subsequent \texttt{ToolMessage} provides the output.  In the
Streamlit app, these are rendered as expandable cards in the right column:

\begin{lstlisting}
for i, tc in enumerate(tool_calls, 1):
    with st.expander(f"{i}. {tc['tool_name']}"):
        st.code(tc["tool_input"], language="json")
        if tc["tool_output"]:
            st.text(tc["tool_output"][:500])
\end{lstlisting}

A statistics row shows the total tool calls, unique tools used, and message
count.

\subsection{Exercise 2 – Critic Agent}

\textbf{Goal:} Automatically validate the generated travel plan with a second
LLM pass.

\textbf{Implementation:} A separate function \texttt{run\_critic\_agent()}
sends the travel plan to the same Ollama model but with a different
\emph{critic} system prompt:

\begin{lstlisting}
CRITIC_PROMPT = """
You are a travel plan critic. Review the plan and provide:
1. Feasibility score (1-10)
2. Budget assessment
3. Activity balance
4. Potential issues
5. Suggestions (2-3 concrete improvements)
"""
\end{lstlisting}

The critic does \emph{not} use tools—it is a pure LLM evaluation.  This
two-agent pattern (generator + critic) is a lightweight form of
\emph{self-reflection} that improves output quality without fine-tuning.

\subsection{Exercise 3 – Budget Constraint}

\textbf{Goal:} Let the user set a maximum budget; the agent must respect it.

\textbf{Implementation:}
\begin{enumerate}[nosep]
  \item The Streamlit sidebar exposes a toggle and a numeric input for the
        budget limit (USD).
  \item When enabled, an instruction is appended to the user prompt:
        \emph{``You MUST keep the total trip cost under \$X.''}
  \item If the plan appears to exceed the budget, the user can click
        \textbf{Re-plan with Tighter Budget}, which re-invokes the agent
        with 90\% of the original limit and an explicit instruction to
        use cheaper alternatives.
\end{enumerate}

% ══════════════════════════════════════════════════════════════════════════════
\section{Streamlit User Interface}
% ══════════════════════════════════════════════════════════════════════════════

The web UI (\texttt{src/agent/app.py}) is divided into:

\begin{description}[nosep]
  \item[Header] Gradient banner with project title.
  \item[Sidebar] Model selector, budget toggle, critic toggle, MCP server
        status reference.
  \item[Main area] Text area with example prompts, a ``Plan My Trip''
        button, stats row, two-column layout (plan + tool calls), critic
        review section, and budget re-planning.
\end{description}

Custom CSS gives the interface a dark-themed, modern aesthetic with
Catppuccin-inspired colours.

% ══════════════════════════════════════════════════════════════════════════════
\section{Running the Project}
% ══════════════════════════════════════════════════════════════════════════════

\begin{enumerate}
  \item \textbf{Install dependencies:}
        \begin{verbatim}
pip install -e .
        \end{verbatim}
  \item \textbf{Start Ollama:}
        \begin{verbatim}
ollama serve
ollama pull llama3.1
        \end{verbatim}
  \item \textbf{Launch MCP servers:}
        \begin{verbatim}
python scripts/run_servers.py
        \end{verbatim}
  \item \textbf{Run the Streamlit app:}
        \begin{verbatim}
streamlit run src/agent/app.py
        \end{verbatim}
\end{enumerate}

% ══════════════════════════════════════════════════════════════════════════════
\section{Discussion}
% ══════════════════════════════════════════════════════════════════════════════

\subsection{Strengths}

\begin{itemize}[nosep]
  \item \textbf{Modularity}: Each MCP server is independent, testable, and
        replaceable (e.g., swap the weather stub for a real API).
  \item \textbf{Transparency}: The tool-call log gives full observability
        into the agent's reasoning.
  \item \textbf{Local execution}: Ollama keeps all data on-device—no API
        keys, no cloud dependency.
  \item \textbf{Self-reflection}: The critic agent adds a lightweight
        quality gate.
\end{itemize}

\subsection{Limitations}

\begin{itemize}[nosep]
  \item \textbf{Simulated data}: The MCP servers return hardcoded data for
        five cities only.
  \item \textbf{Budget enforcement}: The budget is enforced via prompt
        engineering, not programmatically.  The agent may still hallucinate
        numbers.
  \item \textbf{No memory}: Each query is stateless; there is no
        conversation history across requests.
  \item \textbf{Model quality}: Smaller Ollama models may struggle with
        complex multi-tool plans.
\end{itemize}

\subsection{Possible Extensions}

\begin{itemize}[nosep]
  \item Connect MCP servers to real APIs (e.g., OpenWeatherMap, Skyscanner).
  \item Add a \textbf{planner agent} that decomposes complex queries into
        sub-tasks before the executor agent runs.
  \item Implement \textbf{conversation memory} with LangGraph checkpointing.
  \item Add \textbf{unit tests} for each MCP server and integration tests
        for the full agent loop.
  \item Containerise with Docker Compose for one-command deployment.
\end{itemize}

% ══════════════════════════════════════════════════════════════════════════════
\section{Conclusion}
% ══════════════════════════════════════════════════════════════════════════════

This lab demonstrated how to build a fully functional \textbf{agentic AI
system} by combining an LLM (Ollama / LLaMA~3.1), a graph-based agent
framework (LangGraph), and standardised tool servers (MCP).  The three
exercises highlighted key agentic patterns:

\begin{enumerate}[nosep]
  \item \textbf{Tool-augmented generation} – extending the LLM's
        capabilities with external tools.
  \item \textbf{Self-reflection} – using a critic agent to evaluate and
        improve outputs.
  \item \textbf{Constrained planning} – enforcing user-defined limits
        through prompt engineering and iterative re-planning.
\end{enumerate}

The modular architecture (separate MCP servers, clean \texttt{src/} layout,
Streamlit UI) makes the system easy to extend, test, and deploy.

% ══════════════════════════════════════════════════════════════════════════════
% References
% ══════════════════════════════════════════════════════════════════════════════
\begin{thebibliography}{9}

\bibitem{yao2023react}
  S.~Yao, J.~Zhao, D.~Yu, N.~Du, I.~Shafran, K.~Narasimhan, and Y.~Cao,
  ``ReAct: Synergizing Reasoning and Acting in Language Models,''
  in \emph{ICLR}, 2023.

\bibitem{langchain}
  LangChain, ``LangChain Documentation,''
  \url{https://docs.langchain.com/}, 2024.

\bibitem{langgraph}
  LangChain, ``LangGraph: Build Stateful Agents,''
  \url{https://langchain-ai.github.io/langgraph/}, 2024.

\bibitem{mcp}
  Anthropic, ``Model Context Protocol,''
  \url{https://modelcontextprotocol.io/}, 2024.

\bibitem{ollama}
  Ollama, ``Ollama – Run LLMs Locally,''
  \url{https://ollama.ai/}, 2024.

\bibitem{streamlit}
  Streamlit Inc., ``Streamlit Documentation,''
  \url{https://docs.streamlit.io/}, 2024.

\end{thebibliography}

\end{document}
